\chapter{Coronary Artery Disease}
\index{Coronary Artery Disease}
\label{sec:Coronary Artery Disease}

\section{Overview}
Coronary artery disease (CAD), also called coronary heart disease or ischemic heart disease, occurs when plaque narrows or blocks the arteries supplying the heart muscle. Reduced blood flow may cause stable angina, while sudden plaque disruption and clot formation can cause an acute coronary syndrome, including myocardial infarction. CAD may be present without symptoms until a serious event occurs.
\section{Causes}
The usual cause is atherosclerosis. Injury or dysfunction of the arterial lining permits lipid accumulation and inflammation, producing plaques that narrow the coronary arteries. Plaque rupture can trigger platelet activation and thrombosis. Less common causes include coronary spasm, spontaneous coronary artery dissection, and coronary microvascular disease.
\section{Risk Factors}
\begin{itemize}[noitemsep]
  \item Smoking, hypertension, diabetes, and abnormal blood lipids
  \item Increasing age and family history of premature cardiovascular disease
  \item Obesity, physical inactivity, and an unhealthy diet
  \item Chronic kidney disease and inflammatory disorders
  \item Previous stroke, peripheral artery disease, or other atherosclerotic disease
\end{itemize}
\section{Signs and Symptoms}
Typical angina is pressure, tightness, heaviness, or discomfort in the chest brought on by exertion or emotional stress and relieved by rest. Discomfort may spread to the arm, shoulder, neck, jaw, or back. Shortness of breath, fatigue, nausea, or sweating may occur, particularly in older adults, women, and people with diabetes. New, severe, prolonged, or rest pain may indicate a heart attack and requires emergency assessment.
\section{Progression and Stages}
Early atherosclerosis may be silent. As plaque enlarges, exertion can exceed the heart's blood supply and produce stable angina. Plaque rupture may abruptly reduce flow and cause unstable angina or myocardial infarction. Chronic ischemia can contribute to ischemic cardiomyopathy and heart failure.
\section{Complications}
\begin{itemize}[noitemsep]
  \item Myocardial infarction and sudden cardiac arrest
  \item Heart failure and ischemic cardiomyopathy
  \item Ventricular or atrial arrhythmias
  \item Cardiogenic shock
  \item Stroke and peripheral arterial disease as manifestations of systemic atherosclerosis
\end{itemize}
\section{Diagnosis}
Assessment begins with the symptom pattern, medical and family history, physical examination, and cardiovascular risk estimation. Tests may include an electrocardiogram, blood tests including cardiac troponin when acute ischemia is suspected, echocardiography, stress testing, coronary CT angiography, and coronary angiography. The choice depends on symptoms, urgency, pretest probability, and comorbidities.
\section{Prevention}
\begin{itemize}[noitemsep]
  \item Do not smoke; offer cessation support when needed
  \item Follow a heart-healthy eating pattern and exercise regularly
  \item Control blood pressure, diabetes, and lipid abnormalities
  \item Maintain a healthy weight and sleep adequately
  \item Take preventive medicines when prescribed
\end{itemize}
\section{Treatment Options}
Treatment may include lifestyle change, lipid-lowering therapy, antiplatelet therapy when indicated, antianginal medicines, and medicines for blood pressure or diabetes. Acute coronary syndromes require emergency treatment. Revascularization with percutaneous coronary intervention or coronary artery bypass grafting is considered when symptoms or coronary anatomy warrant it. Cardiac rehabilitation supports recovery and long-term risk reduction.
\section{Living with It}
Patients should follow the prescribed medication plan, attend follow-up visits, maintain regular activity appropriate to their condition, and learn the warning signs of myocardial infarction. Sudden chest pressure, shortness of breath, cold sweating, fainting, or pain radiating to the arm or jaw warrants emergency services.
\section{Outlook}
CAD is chronic, but progression and cardiovascular events can often be reduced through smoking cessation, risk-factor control, evidence-based medicines, and revascularization when appropriate. Prognosis depends on the extent of disease, ventricular function, associated conditions, and adherence to treatment.
